\documentclass[a4paper, 12pt, french]{report}

% ========================
% IMPORTATION DES PACKAGES
% ========================

\usepackage[T1]{fontenc}
\usepackage[utf8]{inputenc}
\usepackage[margin=2cm]{geometry}
\usepackage{lmodern}
\usepackage{theorem}
\usepackage{amsfonts, amsmath, amssymb, dsfont}
\usepackage{babel}

\pagestyle{headings}

\ifpdf
\usepackage[pdftex=true, hyperindex=true, colorlinks=false]{hyperref}
\else
\usepackage[pdftex=true, hyperindex=true, colorlinks=true]{hyperref}
\fi

\renewcommand{\familydefault}{\sfdefault}

\theoremstyle{break}
\theoremheaderfont{\scshape}
\theorembodyfont{\upshape}

\newtheorem{defin}{Définition}[section]
\newtheorem{prop}[defin]{Proposition}
\newtheorem{coro}[defin]{Corollaire}
\newtheorem{theo}[defin]{Théorème}

\title{PROBABILITÉ}
\author{}
\date{\today}

\begin{document}

\maketitle

\textbf{PRÉCISION} \\
Ceci n'est pas à priori un book, cette document est un réceuille de note pour comprendre et approfondir mes connaissances en mathématiques.
Ce document est crée à partir du language de programmation latex. Même si c'est un note si vous trouvez que cela peut vous être aider que ce soit dans vous exercice ou juste pour comprendre.
Vous pouvez le consulter et même le télécharger.
Mais je précise que, j'ai les reformuler moi même selon mes propres compréhension.

\tableofcontents

\chapter{ESPACE PROBABILISÉ}

\section{Vocabulaire}

\begin{center}
	\begin{tabular}{|p{8cm}|p{3cm}|}
	\hline
	Vocabulaire & Notation \\ \hline
	Résultat possible & $\omega$ \\
	Tous les résultat possible & $\Omega$ \\
	Évenement A & A \\
	Évenement contraire & $A^{c}$ \\
	Sous ensemble d'évenement & $A \subset \Omega$ \\
	Événement certain & $\Omega$ \\
	Évenement impossible & $\emptyset$ \\
	Évenement A ou B (non exclusif) & $A \cup B$ \\
	Évenement A et B & $A \cap B$ \\
	\hline
\end{tabular}
\end{center}
Prénons un exemple, si on lance un dé à 6 face. On a $\omega = 5$, $\Omega = \{1,2,3,4,5,6\}$, $A = \{1,3,5\}$.\\
On dit que deux évenements sont incompatibles si $A \cap B = \emptyset$.

\section{Probabilité}
$\mathcal{F} = P(\Omega), \; \mathcal{F}$ est ici une partie de $\Omega$ où $\Omega$ est l'ensemble de tous les résultats possibles. $\mathcal{F}$ est appelé un tribus et ($\mathcal{F}, \Omega$) un espace mesurable.

\begin{defin}
	Une mesure de probabilité (ou probabilité) est une application P de $\mathcal{F}$ vers [0, 1] qui verifie les conditions suivantes:
	\begin{enumerate}
		\item P($\Omega$) = 1 et P($\emptyset$) = 0,
		\item Soit $(A_{i})_{i \in \mathbb{N}}$ un collection de famille d'évenement finie ou dénombrable. Si les $A_{i}$ sont deux à deux disjoints, on a $P(\displaystyle\bigcup_{i \in  \mathbb{N}}A_{i}) = \displaystyle\sum_{i \in \mathbb{N}} P(A_{i})$ ($\sigma$-additivité)
	\end{enumerate}
	Le triplet $(\Omega, \mathcal{F}, P)$ est un espace de probabilité. On dit qu'un probabilité d'un évenement A est presque sûr si $P(A) = 1$, il est négligeable si $P(A)=0$.
\end{defin} 

\begin{prop}
	\begin{enumerate}
		\item $P(A^{c}) = 1 - P(A)$,
		\item $P(A \cup B) = P(A) + P(B) - P(A \cap B)$,
		\item Si $A \subset B$, alors $P(A) \leq P(B)$,
		\item Soit $(A_{i})_{i \in \mathbb{N}}$ une famille d'évenement finie ou dénombrable deux à deux disjoints tel que $\displaystyle\sum_{i \in \mathbb{N}} P(A_{1}) = 1$. On a\\ $P(B) = \displaystyle\sum_{i \in \mathbb{N}} P(A_{i} \cap B)$ (\textbf{Formule de décomposition})
	\end{enumerate}
\end{prop}
\textit{Preuve.}
\begin{enumerate}
	\item Montrons que $P(A^{c}) = 1 - P(A)$,\\
	On sait que $A \cup A^{c} = \Omega$, Comme $A$ et $A^{c}$ sont disjoints. D'après $\sigma$-additivité, on a $P(A \cup A^{c}) = P(A) + P(A^{c})$. Or $P(A \cup A^{c})=P(\Omega)=1$. Donc $P(A^{c}) = 1 - P(A)$.
	\item Montrons que $P(A \cup B) = P(A) + P(B) - P(A \cap B)$,\\
		On peut décomposer $A \cup B$ comme suit, $A \cup B = A \cup (B \cap A^{c})$. A et $B \cap A^{c}$ sont disjoints. D'après $\sigma$-additivité, on a $P(A \cup (B \cap A^{c})) = P(A) + P(B \cap A^{c})$. En décomposant aussi B, on a $B = (B \cap A) \cup (B \cup A^{c})$. Par suit, on a $P(B) = P(B \cap A) \cup (B \cup A^{c}) = P(B \cap A) + P(B \cup A^{c})$. On a alors $P(B \cup A^{c}) = P(B) - P(A \cap B)$. D'où $P(A \cup B) = P(A) + P(B) - P(A \cap B)$.
    \item Montrons que si $A \subset B$, alors $P(A) \leq P(B)$,\\
     Supposons que $A \subset B$, considérons deux ensembles deux à deux disjoints A et B \textbackslash A. On a $B = A \cup (B \cap A^{c})$ (car par hypothèse $A \cup B = B$). Par suite $P(B) = P(A) + P(B \cap A^{c}$). Comme $P(B \cap A^{c}) \geq 0$. \\Donc on peut en conclure que $P(B)>P(A)$.
     \item On peut écrire $B = B \cap \Omega =  B \cap (\displaystyle\bigcup_{i \in I} A_{i})$. D'après la distributivité on a $\displaystyle B = \bigcup_{i \in I} (B \cap A_{i})$. 
     Comme les $A_{i}$ sont deux à deux disjoints alors ($B \cap A_{i}$) sont deux à deux disjoints. Donc on a $P(B) = P(\displaystyle\bigcup_{i} (B \cap A_{i}) = \displaystyle\sum_{i \in I} P(B \cap A_{i}))$.\\
     D'où $\displaystyle P(B) = \displaystyle\sum_{i \in I} P(B \cap A_{i})$.  \marginpar{$\square$}
\end{enumerate}

\section{Probabilité sur des ensembles finie ou dénombrable}
\begin{prop}
     Pour tout $A \subset \mathcal{F}$, on a $P(A) = \displaystyle\sum_{\omega \in A} P(\{ \omega\})$.
\end{prop}
\textit{Preuve.}\\
A peut s'écrire sous la forme $A = (A \cap \omega)_{\omega \in \Omega} = (A \cap \{\omega\})_{\omega \in A} \cup (A \cap \{\omega\})_{\omega \not\in A}$. Ces deux ensembles sont deux à deux disjoints. Donc d'après $\sigma$-additivité, on a
$P(A) = P( (A \cap \{\omega\})_{\omega \in A} \cup (A \cap \{\omega\})_{\omega \not\in A}) = P((A \cap \{\omega\})_{\omega \in A}) + P((A \cap \{\omega\})_{\omega \not\in A}) = \displaystyle\sum_{\omega\in A} P(A \cap \{ \omega \}) + \displaystyle\sum_{\omega \not\in A} P(A \cap \{ \omega \}) = \sum_{\omega \in A} P(\{ \omega \})$.\\
D'où $P(A) = \displaystyle\sum_{\omega \in A} P(\{ \omega\})$.

\begin{defin}
        $P(\{ \omega \}) = \frac{1}{|\Omega|}$
\end{defin}

\begin{coro}
       $P(A) = \frac{|A|}{|\Omega|}$
\end{coro}

\section{Dénombrement}
\begin{enumerate}
       \item Le nombre de permutation de $\{1, 2, \cdots, n\}$ dans lui même est $n!$.
       \item Le nombre d'arragement $A_{n}^{k} = \displaystyle\frac{n!}{(n-k)!}$.
       \item Le coefficient binomial $C_{n}^{k} = \displaystyle\frac{n!}{k!(n-k)!}$.
\end{enumerate}

\section{Probabilité conditionnelle et Indépendance}
\begin{defin}
     $P(A | B) = \frac{P(A \cap B)}{P(B)}$
\end{defin}

\begin{prop}
\textbf{Formule de décomposition :} $P(A) = P(A | B)P(B) + P(A | B^{c}) P(B^{c})$ \\
\textbf{Formule de Bayes :} $ P(B | A) = \displaystyle\frac{P(A | B) P(B)}{P(A | B)P(B) + P(A | B^{c}) P(B^{c})}$
\end{prop}

\begin{defin}
On dit que deux événements A et B sont indépendants si $P(A \cap B) = P(A)P(B)$
\end{defin}
Si deux événements A et B sont indépendants, alors la probabilité conditionnelle devient $P(A|B)=P(A)$.

\chapter{VARIABLES ALÉATOIRES}
\section{Lois usuels}
\begin{itemize}
	\item Loi de Bernouilli de paramètres p: \\
	$P(X=0) = 1-p$ et $P(X=1)=p$

	\item Loi uniforme: \\
	$P(X = k)=\displaystyle\frac{1}{N}$.

	\item Loi Binomiale: \\
	$P(X=k)= \left(\begin{array}{c} n \\ k\end{array}\right) p^{k}(1-p)^{n-k}$.

	\item Loi de poison de paramètres $\lambda$: \\
	$P(X=k)=e^{-\lambda} \frac{\lambda^{k}}{k!}$

	\item Loi géometrique de paramètres: $\lambda$: \\
	$P(X=k)=(1-\lambda)^{k-1}\lambda$.
\end{itemize}

\section{Lois marginales}

\section{Indépendence}

\section{Esperance et variance}

\section{Fonction génératrice}

\section{Lois faibles des grands nombres}

\chapter{Variable aleatoire à densité ou continu}

\section{Définition}
$\displaystyle \int_{\mathbb{R}} f(x) dx = 1$.

\section{Lois marginales}
$ \displaystyle f_{X_{1}}(x_{1}) := \int_{\mathbb{R}^{d}} f_{X_{1}, X_{2}, \cdots, X_{n}}(x_{1}, x_{2}, \cdots, x_{n}) dx_{2}dx_{3}\cdots dx_{n}$.

\section{Indépendance}
$\displaystyle f_{X_{1}, X_{2}, \cdots, X_{n}} = f_{X_{1}}(x_{1})f_{X_{2}}(x_{2}) \cdots f_{X_{n}}(x_{n})$.

\section{Esperence et variance}
$\displaystyle E(g(X)) := \int_{\mathbb{R}} g(x)f(x) dx$

$\displaystyle Var(X) := E(X^{2}) - (E(X))^{2}$

\section{Lois usuels}
\begin{itemize}
	\item Loi uniforme $\mathcal{U}([a,b])$: \\ $f(x) := \frac{1}{b-a} \mathds{1}_{[a.b]}(x)$
	\item Loi exponentielle $\epsilon(\lambda)$: \\ $f(x) := \lambda \exp{-\lambda x}$
	\item Loi Gamma $\gamma(\lambda, a)$: \\ $f(x) := \frac{\lambda^{a}}{\Gamma (a)} x^{a-1} \exp{-\lambda x} \mathds{1}_{x>0}$ avec $\Gamma (a) := \int_{0}^{\infty} x^{a-1}\exp^{-\lambda x} dx$.
	\item Loi de beta $\beta (a,b)$: \\ $f(x):= \frac{\Gamma (a+b)}{\Gamma(a) \Gamma{b}} x^{a-1} (1-x)^{b-1}\mathds{1}_{]0,1[}(x)$
	\item Loi normale ou Gaussienne $\mathcal{N}(\mu, \sigma^{2})$: \\ $f(x):= \frac{1}{\sqrt{2\pi}}\sigma \exp{(\frac{\mu - x}{\sigma})^{2}} \mathds{1}_{x>0}(x)$.
\end{itemize}

\chapter{Fonction caracteristique}

\section{Définition}

\section{Propriété}

\section{Indépendance}

\chapter{Convergence en loi}

\end{document}
